\section{Gabarito e resoluções comentadas — Auditoria do Setor Público}

\begin{solucao}{19.01}\textbf{CERTO.} Critério é o padrão de comparação; condição é o estado encontrado. Confundir os dois produz achado mal estruturado.\end{solucao}
\begin{solucao}{19.02}\textbf{ERRADO.} Ceticismo é postura questionadora e avaliação crítica da evidência, não presunção automática de fraude.\end{solucao}
\begin{solucao}{19.03}\textbf{CERTO.} Julgamento profissional aplica competência e experiência sob critérios e normas, devendo ser justificável e documentado.\end{solucao}
\begin{solucao}{19.04}\textbf{CERTO.} Sigilo, fraude, alta administração ou serviço essencial podem tornar um tema material mesmo com baixo valor financeiro.\end{solucao}
\begin{solucao}{19.05}\textbf{CERTO.} Suficiência trata quantidade; apropriação, qualidade — relevância e confiabilidade.\end{solucao}
\begin{solucao}{19.06}\textbf{ERRADO.} Indagação ajuda a compreender, mas controle crítico exige corroboração com registros, observação, reexecução ou outras fontes.\end{solucao}
\begin{solucao}{19.07}\textbf{CERTO.} Teste de controle avalia o controle; procedimento substantivo testa diretamente dados, saldos ou transações.\end{solucao}
\begin{solucao}{19.08}\textbf{ERRADO.} Início de providência não comprova implementação nem efetividade. Monitoramento deve avaliar o que foi realmente implantado e seus efeitos.\end{solucao}
\begin{solucao}{19.09}\textbf{CERTO.} Essa estrutura conecta o que deveria ocorrer, o que ocorreu, por que, com que efeito e qual evidência sustenta a conclusão.\end{solucao}
\begin{solucao}{19.10}\textbf{CERTO.} O auditor deve recomendar o resultado/controle necessário sem assumir a gestão quando existirem alternativas legítimas.\end{solucao}

\begin{solucao}{19.11}\textbf{B.} Economicidade, eficiência e efetividade são o núcleo clássico da auditoria operacional.\end{solucao}
\begin{solucao}{19.12}\textbf{B.} Produzir o mesmo output com menos recursos é melhoria de eficiência.\end{solucao}
\begin{solucao}{19.13}\textbf{CERTO.} Menor preço de aquisição pode gerar custo total maior se a qualidade for inadequada.\end{solucao}
\begin{solucao}{19.14}\textbf{B.} O critério precisa ser norma/política aplicável que defina como a revisão deveria ocorrer.\end{solucao}
\begin{solucao}{19.15}\textbf{CERTO.} Planejamento orienta a coleta para responder questões e sustentar conclusões, evitando trabalho sem propósito.\end{solucao}
\begin{solucao}{19.16}\textbf{B.} Registros e amostras de restauração corroboram a declaração e testam o controle na prática.\end{solucao}
\begin{solucao}{19.17}\textbf{B.} O auditor refaz o cálculo usando dados e fórmula, portanto realiza recálculo.\end{solucao}
\begin{solucao}{19.18}\textbf{CERTO.} Reexecução reproduz atividade/controle; recálculo verifica exatidão matemática ou lógica de cálculo.\end{solucao}
\begin{solucao}{19.19}\textbf{B.} A amostra pode produzir conclusão diferente da população inteira; esse é o risco de amostragem.\end{solucao}
\begin{solucao}{19.20}\textbf{B.} Materialidade qualitativa considera direitos, sigilo, reputação, governança e impacto, não só montante.\end{solucao}
\begin{solucao}{19.21}\textbf{CERTO.} Natureza, época e extensão dos procedimentos são instrumentos do auditor para reduzir risco de detecção.\end{solucao}
\begin{solucao}{19.22}\textbf{A.} Matriz de planejamento liga questão, critério, informação e procedimento, criando rastreabilidade do trabalho.\end{solucao}
\begin{solucao}{19.23}\textbf{A.} A frase é opinativa e ampla. O achado deve indicar critério, condição mensurada e evidência.\end{solucao}
\begin{solucao}{19.24}\textbf{CERTO.} Essa é uma função central da documentação: permitir revisão e compreensão por auditor experiente.\end{solucao}
\begin{solucao}{19.25}\textbf{B.} Publicar política demonstra desenho normativo, não operação. É necessário testar se o processo efetivamente mudou.\end{solucao}

\begin{solucao}{19.26}\textbf{D.} As três assertivas representam elementos fundamentais da auditoria pública.\end{solucao}
\begin{solucao}{19.27}\textbf{CERTO.} No relatório direto, o auditor mede/avalia diretamente; na certificação, examina informação preparada pela parte responsável.\end{solucao}
\begin{solucao}{19.28}\textbf{A.} Há indício de eficiência, mas o efeito pretendido não foi demonstrado, portanto efetividade continua questionável.\end{solucao}
\begin{solucao}{19.29}\textbf{B.} Com RI e RC altos, o auditor precisa reduzir RD para manter RA baixo, aumentando robustez dos procedimentos.\end{solucao}
\begin{solucao}{19.30}\textbf{CERTO.} O modelo orienta raciocínio e planejamento; não precisa de números probabilísticos precisos.\end{solucao}
\begin{solucao}{19.31}\textbf{B.} Screenshot pode ser útil, mas não prova período inteiro, completude nem integridade da origem.\end{solucao}
\begin{solucao}{19.32}\textbf{A.} Se a fonte pode ser alterada pelo auditado, a confiabilidade cai e deve ser reforçada por controles e corroboração.\end{solucao}
\begin{solucao}{19.33}\textbf{CERTO.} Hash demonstra integridade da cópia a partir de um ponto, não veracidade da origem.\end{solucao}
\begin{solucao}{19.34}\textbf{B.} Maior taxa esperada de desvio exige, em geral, mais evidência para sustentar confiança.\end{solucao}
\begin{solucao}{19.35}\textbf{B.} ``Erro humano'' é conclusão prematura. O auditor deve investigar processo, responsabilidades, integração e controles.\end{solucao}
\begin{solucao}{19.36}\textbf{CERTO.} Recomendação baseada em causa errada pode corrigir estoque atual e deixar o mecanismo de recorrência intacto.\end{solucao}
\begin{solucao}{19.37}\textbf{B.} Contraditório é parte do processo de qualidade; evidência nova deve ser analisada e pode mudar o achado.\end{solucao}
\begin{solucao}{19.38}\textbf{B.} A recomendação trata a causa e especifica resultado/control framework sem impor fornecedor.\end{solucao}
\begin{solucao}{19.39}\textbf{CERTO.} Isso preserva a responsabilidade da gestão sobre como implementar a solução.\end{solucao}
\begin{solucao}{19.40}\textbf{A.} Levantamento é adequado para conhecer objeto, estrutura, processos e riscos e preparar trabalhos futuros.\end{solucao}

\begin{solucao}{19.41}\textbf{CERTO.} Três reclamações não sustentam generalização sobre todo o serviço sem critério, universo e procedimentos adicionais.\end{solucao}
\begin{solucao}{19.42}\textbf{B.} O auditor precisa verificar se a revisão realmente ocorreu sobre a população correta e gerou decisões/ações compatíveis.\end{solucao}
\begin{solucao}{19.43}\textbf{B.} Dados de denunciantes possuem relevância qualitativa elevada independentemente de perda financeira.\end{solucao}
\begin{solucao}{19.44}\textbf{CERTO.} Analytics é mecanismo de seleção/evidência complementar. Irregularidade individual ainda exige evidência suficiente e apropriada.\end{solucao}
\begin{solucao}{19.45}\textbf{A.} Sem a versão executada, não se reproduz a análise nem se demonstra exatamente como a evidência foi obtida.\end{solucao}
\begin{solucao}{19.46}\textbf{CERTO.} Código, parâmetros e dados são parte do procedimento quando sustentam a conclusão.\end{solucao}
\begin{solucao}{19.47}\textbf{B.} O estoque foi corrigido, mas a causa preventiva permanece: novos desligamentos continuam gerando contas órfãs.\end{solucao}
\begin{solucao}{19.48}\textbf{A.} Ao gerenciar a implantação, o auditor assume papel gerencial e pode comprometer objetividade em futura avaliação.\end{solucao}
\begin{solucao}{19.49}\textbf{CERTO.} Recomendação pode definir o resultado de controle e deixar a escolha técnica à gestão quando várias alternativas são válidas.\end{solucao}
\begin{solucao}{19.50}\textbf{C.} A formulação conecta condição, causa provável, risco e próximos procedimentos, evitando acusações sem base ou solução simplista.\end{solucao}
